% Secao: KL no talude, analitica contra numerica (agente kl_talude). Sem preambulo; \input no relatorio.
% Fontes: Projects/KLCilindroMika/README.md (secao 4), relatorio/dados/kl_*.csv (kl_truncamento, kl_espectros: todos os
% autovalores de Galerkin; kl_galerkin_talude_*: resumo por M das rodadas do KLTaludeCeron refeitas para o relatorio) e
% Projects/KLCilindroMika/figuras/talude_figuras.py e malhas_campos.py.
\section{Karhunen-Loève no talude: analítica contra numérica}\label{sec:kl:tal}

O \texttt{SlopeRandomField} gera os campos com a forma fechada do retângulo envolvente $R = [0, 25] \times [0, 10]$,
avaliada no talude $\Omega$ (Seção~\ref{kl:form:analitica}). Vargas Ceron et al.\ (2025) obtêm a KL ``from
discretization of the actual soil mass slope'', isto é, numericamente no próprio talude. Esta seção compara as duas, com
a KL exata do talude (Nyström, Seção~\ref{kl:form:nystrom}) como referência. A KL numérica é a de Galerkin do NeoPZ
(Seção~\ref{kl:form:galerkin}, programa \texttt{KLTaludeCeron}) em três malhas de triângulos do talude
(\texttt{ceron\_h5\_a45}, \texttt{\_h025} e \texttt{\_fina}, com $h$ de 0,375, 0,25 e \SI{0.1875}{m} até \SI{1}{m} da
superfície), P1 e P2, com a regra do elemento (ordem $2P$) e com a regra de ordem 10 (25 pontos). Os números são os do
README do \texttt{KLCilindroMika} (Seção~4), conferidos de forma independente em numpy/scipy; as figuras usam as
rodadas do \texttt{KLTaludeCeron} refeitas para este relatório (script \texttt{talude\_figuras.py}), que repetem os
valores do README.

\subsection{Autovalores}

Os 10 maiores autovalores da KL exata do talude são 48,810; 26,730; 16,911; 10,824; 9,572; 6,824; 4,719; 4,519; 3,829 e
3,349; os da forma fechada no retângulo, 57,155; 36,231; 21,400; 13,120; 13,046; 8,570; 8,270; 5,939; 4,885 e 4,511
(README, Seção~4). Os dois conjuntos não são comparáveis termo a termo: a soma de todos os autovalores é a área do
domínio (187,5 no talude e 250 no retângulo), e os modos do retângulo não são ortogonais no talude. Os modos do talude
(Figura~\ref{fig:kl:modostal}) se parecem com os do retângulo (Figura~\ref{fig:kl:modosret}) longe da face e se
deformam perto dela: não são produtos $f_i(x) g_j(y)$. A
Figura~\ref{fig:kl:espectros}(a) mostra os espectros. Os autovalores de Galerkin acompanham os do talude nos primeiros
modos e, perto do número de equações $N$, caem rapidamente: a malha só tem $N$ autovalores, e a sua soma (o traço
numérico) fica abaixo da área (0,9713 da área em P2 \texttt{\_h025} com a regra de ordem 10). A
Figura~\ref{fig:kl:espectros}(b) compara os autovalores de Galerkin com os de Nyström. Nos 10 primeiros modos, com a
regra de ordem 10, a diferença relativa vai de $-2{,}6 \times 10^{-5}$ a $+3{,}0 \times 10^{-5}$ em P2
\texttt{\_h025} e chega a $-1{,}6 \times 10^{-3}$ em P1 \texttt{\_fina}; com a regra do elemento, os autovalores
ficam acima dos exatos (de $+4{,}3 \times 10^{-5}$ a $+8{,}2 \times 10^{-4}$ em P2 \texttt{\_h025}). A diferença
cresce com o índice do modo: de $k = 50$ a 100, nos seis casos da figura, a sua mediana vai de 5\,\% (P2
\texttt{\_h025}, regra do elemento) a 18\,\% (P1 \texttt{\_fina} e P2 \texttt{ceron\_h5\_a45}, regra de ordem 10),
acima da incerteza da referência (P1 \texttt{ceron\_h5\_a45} e P1 \texttt{\_h025} com a regra do elemento, fora da
figura, chegam a 34 e 23\,\%).
Esses modos variam depressa em $y$ ($L_y = \SI{2}{m}$; no retângulo, o índice 1D em $y$ chega a 477 nos 2000 maiores
termos) e não são resolvidos pelos elementos de 0,75 a \SI{1.5}{m} longe da superfície.

\begin{figure}[H]
  \centering
  \includegraphics[width=\textwidth]{kl_modos_talude.png}
  \caption{Os 9 primeiros modos da KL exata do talude (Nyström, grade $n = 160$, modos nos nós de uma grade de
    \SI{0.25}{m}; $\lambda_k$ extrapolados). Dados \texttt{dados/kl\_nystrom\_modos.csv}; scripts \texttt{nystrom.py} e
    \texttt{talude\_figuras.py}.}
  \label{fig:kl:modostal}
\end{figure}

\begin{figure}[H]
  \centering
  \includegraphics[width=\textwidth]{kl_espectros.pdf}
  \caption{(a) Autovalores $\lambda_k$: forma fechada no retângulo $R$, KL exata do talude (Nyström, 210 autovalores,
    extrapolação de Richardson das grades $n = 80$, 160 e 320) e KL de Galerkin no talude. (b) Diferença relativa dos
    autovalores de Galerkin contra os de Nyström, $k \le 100$; a faixa cinza é a diferença entre as extrapolações com
    $n = 40$, 80, 160 e com $n = 80$, 160, 320 (incerteza da referência): diferenças abaixo dela não são resolvidas.
    Script \texttt{talude\_figuras.py}.}
  \label{fig:kl:espectros}
\end{figure}

\subsection{Erro de truncamento}

A Figura~\ref{fig:kl:truncamento} e a Tabela~\ref{tab:kl:truncamento} dão $\varepsilon_M$ (Eq.~\ref{eq:kl:erro}).
\begin{itemize}
  \item Com $M = 50$, a KL numérica no talude tem $\varepsilon_M$ de 0,099 a 0,104 (P1 \texttt{\_fina} e P2, com as duas
    regras), contra 0,1150 da forma fechada do retângulo; a KL exata do talude dá 0,0993, e a forma fechada precisa de
    63 termos para esse erro.
  \item Com $M = 200$, a KL numérica ainda muda com a malha (P2 com a regra de ordem 10: 0,0547 com 2320 equações e
    0,0460 com 5085); a forma fechada dá 0,0444 e a KL exata do talude, 0,0375 (a forma fechada precisa de 253 a 254
    termos para esse erro). Nenhuma das malhas chega a 0,0375 com $M = 200$ (o menor valor é 0,0421).
  \item Com $M = 1000$ e 2000, a KL numérica para no erro da malha: $1 - \text{traço}/|\Omega| = 0{,}0287$ na malha de 5085
    equações com a regra de ordem 10. A forma fechada continua caindo: 0,0135 com 1000 termos, 0,0079 com 2000 (o $M$ do
    programa) e 0,0013 com 20\,000 (\texttt{talude\_figuras.py}).
  \item A regra do elemento dá erros menores que a de 25 pontos na mesma malha porque aumenta os autovalores
    (Seção~\ref{kl:form:galerkin}): com ela, $\varepsilon_{50} = 0{,}0990$ em P2 \texttt{\_h025}, abaixo do mínimo
    teórico (0,0993, o da KL exata do talude). Esses valores medem a soma dos autovalores numéricos, não a qualidade dos
    modos.
\end{itemize}

\begin{figure}[H]
  \centering
  \includegraphics[width=\textwidth]{kl_truncamento.pdf}
  \caption{Erro médio quadrático de truncamento no talude, $\varepsilon_M$, contra o número de termos: forma fechada do
    retângulo restrita ao talude (integrais exatas de cada modo no talude, até 20\,000 termos), KL exata do talude
    (Nyström extrapolado, até 210 termos) e KL de Galerkin no talude ($1 - \sum_{k \le M} \lambda_k^h / 187{,}5$, com
    todos os autovalores de cada malha), (a) com a regra do elemento e (b) com a de ordem 10. Script
    \texttt{talude\_figuras.py}.}
  \label{fig:kl:truncamento}
\end{figure}

% Tabela (agente kl_talude): erro medio quadratico de truncamento no talude. Fonte: Projects/KLCilindroMika/README.md,
% secao 4 (tabelas "Erro medio quadratico de truncamento no talude" e "Termos da forma fechada com o mesmo erro").
\begin{table}[H]
  \centering
  \caption{Erro médio quadrático de truncamento no talude, $\varepsilon_M$ (Eq.~\ref{eq:kl:erro}), e número de termos
    $M_R$ da forma fechada do retângulo com o mesmo $\varepsilon$ que a KL do talude com $M$ termos. Galerkin: KL numérica
    do NeoPZ no talude (malhas \texttt{ceron\_h5\_a45}: $h = 0{,}375$ e \SI{1.5}{m} até \SI{1}{m} e a partir de
    \SI{6}{m} da superfície; \texttt{\_h025}: 0,25 e \SI{1}{m}; \texttt{\_fina}: 0,1875 e \SI{0.75}{m}); regra: ordem/pontos
    da regra do triângulo (2/3 e 4/6: a regra do elemento, ordem $2P$); ``todos'': $1 - \text{traço}/|\Omega|$ com todos os
    autovalores numéricos. Na KL numérica, $\varepsilon_M$ sai dos autovalores numéricos. Nyström: KL exata do talude,
    extrapolada. Fonte: README do \texttt{KLCilindroMika}, Seção~4.}
  \label{tab:kl:truncamento}
  \footnotesize
  \setlength{\tabcolsep}{3.6pt}
  \begin{tabular}{llcrcccccccccc}
    \toprule
    KL & malha & $P$ & eq. & regra & todos & \multicolumn{4}{c}{$\varepsilon_M$} & \multicolumn{4}{c}{$M_R$} \\
    \cmidrule(lr){7-10}\cmidrule(lr){11-14}
     & & & & & & $M = 50$ & 200 & 1000 & 2000 & 50 & 200 & 1000 & 2000 \\
    \midrule
    Galerkin & \texttt{a45}  & 1 &  609 & 2/3   & 0,0715 & 0,1149 & 0,0767 & --     & --     & 51 &  92 & --  & --  \\
    Galerkin & \texttt{h025} & 1 & 1314 & 2/3   & 0,0500 & 0,1060 & 0,0608 & 0,0502 & --     & 57 & 129 & 169 & --  \\
    Galerkin & \texttt{fina} & 1 & 2301 & 2/3   & 0,0382 & 0,1021 & 0,0524 & 0,0390 & 0,0382 & 60 & 159 & 240 & 247 \\
    Galerkin & \texttt{a45}  & 2 & 2320 & 4/6   & 0,0273 & 0,0995 & 0,0463 & 0,0286 & 0,0273 & 63 & 189 & 368 & 391 \\
    Galerkin & \texttt{h025} & 2 & 5085 & 4/6   & 0,0190 & 0,0990 & 0,0421 & 0,0224 & 0,0198 & 63 & 216 & 512 & 604 \\
    Galerkin & \texttt{fina} & 1 & 2301 & 10/25 & 0,0444 & 0,1039 & 0,0560 & 0,0449 & 0,0444 & 59 & 144 & 198 & 200 \\
    Galerkin & \texttt{a45}  & 2 & 2320 & 10/25 & 0,0419 & 0,1030 & 0,0547 & 0,0426 & 0,0419 & 60 & 150 & 213 & 217 \\
    Galerkin & \texttt{h025} & 2 & 5085 & 10/25 & 0,0287 & 0,1002 & 0,0460 & 0,0305 & 0,0291 & 62 & 191 & 337 & 360 \\
    \midrule
    \multicolumn{5}{l}{forma fechada do retângulo $R$, restrita a $\Omega$} & -- & 0,1150 & 0,0444 & 0,0135 & 0,0079 & 50 & 200 & 1000 & 2000 \\
    \multicolumn{5}{l}{exata do talude (Nyström)} & 0 & 0,0993 & 0,0375 & -- & -- & 63 & 253--254 & -- & -- \\
    \bottomrule
  \end{tabular}
\end{table}


\subsection{Variância pontual e covariância}

A Figura~\ref{fig:kl:variancia} e a Tabela~\ref{tab:kl:variancia} mostram como as duas KL distribuem o erro no talude.
A variância $\sigma_M^2$ da forma fechada fica abaixo de 1 em todos os pontos, é suave e o seu intervalo encolhe com $M$
(0,884 a 0,959 com 200 termos; 0,977 a 0,992 com 2000). A da KL numérica passa de 1 e tem a marca dos elementos: com
$M = 2000$, vai de 0,915 a 1,111 na malha de 5085 equações, e de 0,848 a 1,165 em P1 \texttt{\_fina} com a regra do elemento;
com todos os modos, $\sigma^2$ é a diagonal do núcleo projetado na malha, não 1. Os erros da covariância da forma fechada
caem com $M$ (pares próximos, erro médio: $2{,}5 \times 10^{-2}$ com 50 termos e $2{,}8 \times 10^{-4}$ com 2000); os da
KL numérica quase não mudam de $M = 200$ a 2000 ($8{,}5 \times 10^{-3}$ e $7{,}7 \times 10^{-3}$). Com 50 e 200
termos, a KL exata do talude tem os quatro erros da covariância menores que os da forma fechada (com $M = 200$, pares
próximos, erro médio de $3{,}8 \times 10^{-3}$ contra $4{,}9 \times 10^{-3}$), e, com $M = 200$, a KL numérica da malha de 5085 equações tem
erros maiores que os das duas ($8{,}5 \times 10^{-3}$).

\begin{figure}[H]
  \centering
  \includegraphics[width=\textwidth]{kl_variancia.png}
  \caption{Variância truncada $\sigma_M^2(\mathbf{x})$ com $M = 200$ e 2000: forma fechada do retângulo restrita ao
    talude (a, c) e KL de Galerkin no talude, P2, \texttt{ceron\_h5\_a45\_h025}, regra de ordem 10 (b, d), na rede de 15
    pontos de cada triângulo; escala divergente centrada em 1 e, no título, o mínimo e o máximo. Script
    \texttt{talude\_figuras.py}.}
  \label{fig:kl:variancia}
\end{figure}

% Tabela (agente kl_talude): variancia truncada e erros da covariancia truncada no talude. Fonte:
% Projects/KLCilindroMika/README.md, secao 4 (KL numerica P2 ceron_h5_a45_h025, ordem 10; Nystrom n = 160/320).
\begin{table}[H]
  \centering
  \caption{Variância truncada $\sigma_M^2$ (mínimo e máximo na rede de 15 pontos de cada triângulo; Nyström: nos nós da
    grade $n = 160$) e erro da covariância truncada $C_M$ contra o núcleo (máximo e médio em módulo) em 500 pares
    uniformes e 500 pares próximos no talude. Galerkin: P2, \texttt{ceron\_h5\_a45\_h025}, 5085 equações, regra de ordem
    10; a forma fechada é avaliada nos mesmos pontos. Fonte: README do \texttt{KLCilindroMika}, Seção~4.}
  \label{tab:kl:variancia}
  \small
  \begin{tabular}{rlcccccc}
    \toprule
     & & \multicolumn{2}{c}{$\sigma_M^2$} & \multicolumn{2}{c}{pares uniformes} & \multicolumn{2}{c}{pares próximos} \\
    \cmidrule(lr){3-4}\cmidrule(lr){5-6}\cmidrule(lr){7-8}
    $M$ & KL & mín. & máx. & erro máx. & erro médio & erro máx. & erro médio \\
    \midrule
      50 & Galerkin no talude          & 0,769 & 0,992 & \num{9.2e-2} & \num{3.6e-3} & \num{1.1e-1} & \num{2.0e-2} \\
      50 & forma fechada de $R$         & 0,731 & 0,895 & \num{1.0e-1} & \num{4.2e-3} & \num{1.4e-1} & \num{2.5e-2} \\
      50 & exata do talude (Nyström)    & 0,763 & 0,917 & \num{9.1e-2} & \num{3.4e-3} & \num{1.1e-1} & \num{1.9e-2} \\
    \midrule
     200 & Galerkin no talude          & 0,896 & 1,069 & \num{4.8e-2} & \num{1.0e-3} & \num{5.4e-2} & \num{8.5e-3} \\
     200 & forma fechada de $R$         & 0,884 & 0,959 & \num{3.7e-2} & \num{8.0e-4} & \num{4.2e-2} & \num{4.9e-3} \\
     200 & exata do talude (Nyström)    & 0,902 & 0,968 & \num{3.0e-2} & \num{6.5e-4} & \num{3.5e-2} & \num{3.8e-3} \\
    \midrule
    1000 & Galerkin no talude          & 0,915 & 1,105 & \num{3.6e-2} & \num{7.6e-4} & \num{5.5e-2} & \num{7.7e-3} \\
    1000 & forma fechada de $R$         & 0,962 & 0,987 & \num{8.6e-3} & \num{1.0e-4} & \num{1.0e-2} & \num{6.9e-4} \\
    \midrule
    2000 & Galerkin no talude          & 0,915 & 1,111 & \num{3.5e-2} & \num{7.5e-4} & \num{5.5e-2} & \num{7.7e-3} \\
    2000 & forma fechada de $R$         & 0,977 & 0,992 & \num{4.3e-3} & \num{3.7e-5} & \num{4.3e-3} & \num{2.8e-4} \\
    \bottomrule
  \end{tabular}
\end{table}


\subsection{Custo}

A forma fechada com 20\,000 termos e as integrais de cada modo no talude levou 1,3~s (README, Seção~4). A KL de Galerkin
tem memória $O(N^2)$ e solução $O(N^3)$: na malha de 5085 equações, 260~s e 650~MB no README e 216~s e 653~MB nas
rodadas refeitas para este relatório (Tabela~\ref{tab:kl:custo}), quase todo o tempo no problema de autovalores denso; com a regra
de ordem 10 em P1, a montagem de $\mathbf{C}$ pesa mais (62~s de 77~s na malha \texttt{\_fina}). O Nyström de
referência ($n = 160$, 193\,041 nós) levou 81~s e 1,4~GB para 210 autopares.

% Tabela (agente kl_talude): custo das formulacoes no talude. Fontes: Galerkin, dados/kl_galerkin_talude_regra_elemento.csv e
% dados/kl_galerkin_talude_ordem10.csv (rodadas refeitas para o relatorio, KLTaludeCeron, 1 thread);
% Nystrom, dados/kl_nystrom_custo.csv; forma fechada, README do KLCilindroMika, secao 4.
\begin{table}[H]
  \centering
  \caption{Custo de cada formulação no talude (uma thread). Galerkin: KL numérica do NeoPZ (\texttt{KLTaludeCeron}), tempo
    de montagem de $\mathbf{C}$, do problema de autovalores com autovetores e total, e memória máxima, nas rodadas refeitas
    para este relatório; os tempos dependem da carga da máquina (o README registra 260~s e 650~MB para a malha de 5085
    equações com a regra de ordem 10).
    Nyström: $n = 160$, 210 autopares. Forma fechada: 20\,000 termos e as integrais de cada modo no talude.}
  \label{tab:kl:custo}
  \footnotesize
  \setlength{\tabcolsep}{4pt}
  \begin{tabular}{llcrrrrr}
    \toprule
    KL & malha & $P$ / regra & incógnitas & $\mathbf{C}$ (s) & autopares (s) & total (s) & memória (MB) \\
    \midrule
    Galerkin & \texttt{ceron\_h5\_a45}      & 1 / 2  &    609 &  0,11 &   0,22 &   0,71 &  61 \\
    Galerkin & \texttt{ceron\_h5\_a45\_h025} & 1 / 2  &   1314 &  0,58 &   1,99 &   3,36 &  91 \\
    Galerkin & \texttt{ceron\_h5\_a45\_fina} & 1 / 2  &   2301 &  2,39 &  12,03 &  16,10 & 206 \\
    Galerkin & \texttt{ceron\_h5\_a45}      & 2 / 4  &   2320 &  0,46 &  12,24 &  13,33 & 209 \\
    Galerkin & \texttt{ceron\_h5\_a45\_h025} & 2 / 4  &   5085 &  2,96 & 199,83 & 204,06 & 646 \\
    Galerkin & \texttt{ceron\_h5\_a45\_fina} & 1 / 10 &   2301 & 62,15 &  10,84 &  76,67 & 204 \\
    Galerkin & \texttt{ceron\_h5\_a45}      & 2 / 10 &   2320 &  4,43 &  12,10 &  17,55 & 204 \\
    Galerkin & \texttt{ceron\_h5\_a45\_h025} & 2 / 10 &   5085 & 22,44 & 190,99 & 215,57 & 653 \\
    \midrule
    Nyström  & grade $h = \SI{3.125}{cm}$    & --     & 193\,041 & --  &  77,3  &  80,9  & 1414 \\
    forma fechada & --                       & --     & --     & --    & --     &   1,3  & -- \\
    \bottomrule
  \end{tabular}
\end{table}


\subsection{O mesmo campo em duas malhas}

A Figura~\ref{fig:kl:malha} mostra o campo $c'$ da mesma realização (\texttt{SEMENTE} padrão, \texttt{REAL=0}) em duas
malhas: a \texttt{ceron\_h5\_a45\_foco375} (720 triângulos, 2160 pontos de integração) e a malha refinada da realização 0
do método de duas etapas (2544 triângulos, 7632 pontos; $h = \SI{0.125}{m}$ na banda, \SI{0.375}{m} no foco e até
\SI{3}{m} longe do talude). O programa grava os campos nos pontos de integração (\texttt{SO\_CAMPOS=1}), e o mesmo campo
é refeito em Python (\texttt{klfechada.py}: forma fechada, os $\xi$ do gerador do programa e a transformação lognormal)
nos pontos de cada malha (\texttt{dados/kl\_malha\_resumo.csv}):
\begin{itemize}
  \item KL analítica: a maior diferença relativa entre o programa e a função é $3{,}4 \times 10^{-6}$ na
    \texttt{foco375} e $4{,}0 \times 10^{-6}$ na refinada em $c'$ ($1{,}1$ e $1{,}3 \times 10^{-6}$ em $\phi'$; $7{,}7$ e $7{,}3 \times 10^{-6}$ em $k$; os 7 algarismos do arquivo); nos 6
    pontos de integração que as duas malhas têm em comum, $c'$ e $\phi'$ são iguais (diferença 0); e o campo da rodada
    só de campos é o mesmo da simulação completa na mesma malha (diferença 0). As duas malhas veem a mesma função.
  \item KL numérica (contraste): uma KL de Galerkin P1 em numpy (\texttt{galerkin\_p1.py}, $M = 200$) em cada malha, com
    os mesmos $\xi$ de $c'$. Os autovalores e o erro de truncamento quase não mudam ($\lambda_1 = 48{,}825$ e $48{,}830$;
    $\varepsilon_{200} = 0{,}128$ e $0{,}127$), mas o campo muda: nos pontos da malha refinada, a diferença RMS relativa
    entre os dois campos é 0,66 (correlação 0,63), em parte porque o sinal de cada autovetor é arbitrário (97 dos 200
    produtos internos entre modos de mesmo índice são negativos); alinhando o sinal de cada modo (o que só é possível
    tendo as duas malhas), ainda sobra 0,29 (correlação 0,85), porque os modos de mesmo índice deixam de coincidir: $|\int f_k^{\mathrm{foco375}} f_k^{\mathrm{ref}}| \ge 0{,}80$ até $k = 20$,
    0,58 em $k = 40$ e, a partir de $k \approx 50$, valores espalhados entre 0 e 0,87.
\end{itemize}

\begin{figure}[tbp]
  \centering
  \includegraphics[width=\textwidth]{kl_malha_independencia.png}
  \caption{O mesmo campo em duas malhas. (a, b) $c'$ da KL analítica do programa ($M = 2000$) na malha
    \texttt{foco375} e na refinada da realização 0 (média dos pontos de integração de cada triângulo). (c, d) $c'$ de uma
    KL numérica P1 em cada malha ($M = 200$), com os mesmos $\xi$. (e) Programa contra a forma fechada em Python nos
    pontos das duas malhas. (f) Os dois campos numéricos nos mesmos pontos, com o sinal do solver e com o sinal de cada
    modo alinhado. (g) Produto interno dos modos numéricos de mesmo índice das duas malhas. Dados
    \texttt{dados/kl\_malha\_*.csv}; script \texttt{malhas\_campos.py}.}
  \label{fig:kl:malha}
\end{figure}

\subsection{Analítica ou numérica no talude? E a malha?}\label{kl:tal:resposta}

A observação foi: ``se ela [a forma fechada] serve para o talude, não faz sentido usar a numérica, e nem importa o
nível de refinamento da malha para a qualidade dos campos estocásticos''. Os resultados acima a confirmam para este
problema, com as ressalvas abaixo.

\paragraph{A forma fechada serve para o talude.}
Um campo definido no retângulo com covariância $C$ tem, restrito ao talude, exatamente a covariância $C$ no talude: a
covariância é uma propriedade de pares de pontos, e restringir o domínio não a muda. A KL do retângulo, avaliada no
talude, converge para o campo-alvo com $M \to \infty$ ($\varepsilon_M = 0{,}0079$ com 2000 termos e 0,0013 com 20\,000),
e custa 1,3~s contra 216 a 260~s e 650~MB da KL numérica de 5085 equações. Com $M = 2000$, o seu erro de truncamento
(0,0079) é menor que o que a KL numérica atingiu em qualquer das malhas testadas, mesmo com todos os modos (0,0287 na malha
de 5085 equações com a regra de ordem 10; os 0,0190 da regra do elemento vêm de autovalores inflados pela quadratura): a
numérica fica limitada pela malha, e baixar esse limite exige malhas mais finas, com memória $O(N^2)$ e tempo $O(N^3)$.
Os erros da covariância também favorecem a forma fechada já com $M = 200$ (os quatro erros da
Tabela~\ref{tab:kl:variancia} são menores que os da KL numérica da malha de 5085 equações); com $M = 1000$ os erros da
numérica são de 4 a 11 vezes os da forma fechada, e com $M = 2000$, de 8 a 28 vezes; só com $M = 50$ a KL numérica é
melhor nesses erros.

O que a forma fechada perde é a eficiência por termo: a KL do próprio talude é ótima para cada $M$, e o retângulo precisa
de 63 termos para o erro de 50 termos da KL exata do talude e de 253 a 254 para o de 200 (26 a 27\,\% a mais). Isso
só pesa quando o número de variáveis aleatórias é o custo: FORM/SORM, caos polinomial, colocação estocástica,
metamodelos. Em Monte Carlo, como aqui, o custo de cada realização é a solução por elementos finitos; somar 2000
termos nos pontos de integração (cerca de $5 \times 10^7$ multiplicações para os três campos nos 7632 pontos da malha
refinada) é desprezível diante dela. Neste problema (núcleo exponencial separável, talude dentro de um
retângulo, Monte Carlo), usar a KL numérica não traz ganho; nas malhas testadas, ela é menos precisa com $M = 2000$.

\paragraph{Quando a KL numérica continua necessária.}
\begin{itemize}
  \item Núcleos sem forma fechada no domínio, como o exponencial isotrópico $\exp(-r/L)$ do cilindro de Mika et al.\
    (Seção~\ref{sec:kl:cil}), o de Matérn, ou campos não estacionários.
  \item Domínios muito menores que a sua caixa envolvente, em que a perda por termo tende a crescer, porque parte da
    variância de cada modo do retângulo cai fora do domínio (aqui, $|\Omega|/|R| = 0{,}75$ e a perda é de 26 a 27\,\%),
    e domínios não convexos em que a correlação deve ser medida dentro do domínio (por exemplo, através de uma escavação
    ou de uma trinca), onde o núcleo euclidiano restrito não é o modelo pretendido.
  \item Quando o número de variáveis aleatórias precisa ser o menor possível (os métodos acima).
\end{itemize}

\paragraph{A malha e a qualidade dos campos.}
Com a KL analítica, o campo é uma função em forma fechada de $\mathbf{x}$, avaliada em cada ponto de integração: a sua
qualidade ($\varepsilon_M$, $\sigma_M^2$, $C_M$) não depende da malha, e os mesmos $\xi$ dão o mesmo campo em qualquer
malha com o mesmo retângulo envolvente (Figura~\ref{fig:kl:malha}: diferença no nível dos algarismos gravados, e valores
iguais nos pontos comuns). É essa propriedade que torna consistente o refinamento por realização: a malha refinada
amostra a mesma realização da malha grossa. O retângulo é a caixa dos nós; uma malha com outra caixa daria outra
expansão (no talude, todas as malhas têm a caixa $[0, 25] \times [0, 10]$). Com a KL numérica, as autofunções, e com
elas o campo para os mesmos $\xi$, mudam com a malha (Figura~\ref{fig:kl:malha}c, d, f, g): refinar a malha durante
uma realização mudaria a própria realização, e um esquema de refinamento teria de calcular a KL uma vez, numa malha
fixa, e interpolar o campo nas malhas de análise.

\paragraph{O que a malha ainda controla.}
\begin{itemize}
  \item A solução por elementos finitos: o fator de segurança e a banda de cisalhamento, que são o objeto do refinamento
    (Seções~\ref{sec:ref:metodo} e~\ref{sec:ref:niveis}).
  \item Como o modelo mecânico resolve a variabilidade: as propriedades entram nos pontos de integração, e um elemento
    de tamanho $h$ só amostra o campo nos seus pontos. Na banda, $h = 0{,}125$ a \SI{0.375}{m}, bem menor que
    $L_y = \SI{2}{m}$; longe do talude, os elementos chegam a \SI{3}{m} (\texttt{hMax} da malha refinada), maiores que
    $L_y$, e não resolvem as flutuações verticais, mas ali o material fica no regime elástico e pesa pouco na ruptura
    (efeito não quantificado aqui) --- mais uma razão para refinar onde a banda passa.
  \item Nada na malha corrige o déficit de variância da truncação: com $M = 2000$, $\sigma_M^2$ vai de 0,977 a 0,992 no
    talude (média 0,9921), isto é, o desvio padrão de $Z_M$ fica 0,4 a 1,2\,\% abaixo do alvo, em qualquer malha. O
    programa não reescala o campo (\texttt{KL\_variancia} registra a fração retida no retângulo, 0,9920). Se esse déficit
    importar, há duas saídas: aumentar $M$ (o custo por ponto cresce linearmente com $M$) ou dividir $Z_M(\mathbf{x})$
    por $\sigma_M(\mathbf{x})$, que o programa já calcula em cada ponto (\texttt{PreparaPontos}), o que torna exatas as
    marginais lognormais ao custo de uma pequena alteração da correlação.
\end{itemize}
